\documentclass[11pt]{article}

\usepackage[margin=1in]{geometry}
\usepackage{amsmath,amssymb,amsthm}
\usepackage{algorithm}
\usepackage{algpseudocode}
\usepackage{enumitem}
\usepackage{hyperref}
\usepackage{xcolor}
\usepackage{tikz}
\usepackage{booktabs}

\hypersetup{colorlinks=true, linkcolor=blue, urlcolor=blue, citecolor=blue}

\newtheorem{theorem}{Theorem}[section]
\newtheorem{lemma}[theorem]{Lemma}
\newtheorem{proposition}[theorem]{Proposition}
\newtheorem{corollary}[theorem]{Corollary}
\newtheorem{definition}[theorem]{Definition}
\newtheorem{example}[theorem]{Example}

\title{TOAE209/TOAH209 -- Design and Analysis of Algorithms\\[4pt]
\large Week 5: Greedy Algorithms II: Shortest Paths and Minimum Spanning Trees}
\author{Foreign Trade University}
\date{}

\begin{document}
\maketitle
\tableofcontents

\section{Introduction}

Last week introduced the greedy paradigm and two proof techniques --
\emph{greedy stays ahead} and the \emph{exchange argument} -- through scheduling
problems on intervals and jobs. This week we apply the same toolkit to two of
the most important problems in all of algorithms, both defined on
\textbf{weighted graphs}:

\begin{itemize}
    \item the \textbf{single-source shortest path} problem, solved by
    \emph{Dijkstra's algorithm} (a greedy ``stays ahead'' argument); and
    \item the \textbf{minimum spanning tree} (MST) problem, solved by
    \emph{Kruskal's} and \emph{Prim's} algorithms (both justified by a single
    structural fact, the \emph{cut property}).
\end{itemize}

Both algorithms are greedy: they build a solution one vertex or one edge at a
time, always taking the locally cheapest extension, and never reconsidering.
As always, the implementation is short; the content is in the correctness
proof. A recurring sub-theme is that Dijkstra's greedy choice is correct
\emph{only} when all edge weights are non-negative -- a single negative edge can
break it -- which motivates the Bellman--Ford algorithm as a more robust (but
slower) alternative.

Throughout, a graph $G = (V, E)$ has $n = |V|$ vertices and $m = |E|$ edges.
Each edge $e = (u,v)$ carries a weight (or length, or cost) $w(e) = w(u,v) \ge
0$ unless stated otherwise. We represent graphs as \emph{adjacency dicts}: a map
from each vertex to a dict of (neighbour, weight) pairs. The shared helper file
\texttt{starter\_code.py} provides this representation, a seeded random
weighted-graph generator, and a \texttt{UnionFind} data structure; Problem 1
asks you to build and inspect such graphs.

\section{The Single-Source Shortest Path Problem}

\begin{definition}[Single-Source Shortest Paths]
Given a directed or undirected graph $G = (V,E)$ with non-negative edge weights
$w(e) \ge 0$ and a designated \emph{source} vertex $s \in V$, the
\emph{single-source shortest path problem} asks, for every vertex $v \in V$, for
the minimum total weight $\mathrm{dist}(v)$ of any path from $s$ to $v$ (with
$\mathrm{dist}(v) = \infty$ if no path exists). The total weight of a path is the
sum of the weights of its edges.
\end{definition}

This models a road network (weights = travel times), a communication network
(weights = latencies), or any setting where we want the cheapest route from one
origin to all destinations.

\subsection{Dijkstra's algorithm}

Dijkstra's algorithm maintains a set $S$ of vertices whose shortest distance is
already known (``finalized''), and a tentative distance $d[v]$ for every vertex.
Initially $S = \emptyset$, $d[s]=0$, and $d[v]=\infty$ for $v \ne s$. It
repeatedly selects the vertex $u \notin S$ with the smallest tentative distance,
adds it to $S$ (finalizing it), and \emph{relaxes} each outgoing edge $(u,v)$:
if $d[u] + w(u,v) < d[v]$, it lowers $d[v]$.

\begin{algorithm}[h]
\caption{Dijkstra's Single-Source Shortest Paths}
\begin{algorithmic}[1]
\State $d[s] \gets 0$; \quad $d[v] \gets \infty$ for all $v \ne s$
\State $S \gets \emptyset$
\State Insert all vertices into a min-priority queue $Q$ keyed by $d[\cdot]$
\While{$Q \ne \emptyset$}
    \State $u \gets \textsc{ExtractMin}(Q)$ \Comment{smallest tentative distance}
    \State $S \gets S \cup \{u\}$ \Comment{$d[u]$ is now final}
    \For{each edge $(u,v) \in E$ with $v \notin S$}
        \If{$d[u] + w(u,v) < d[v]$}
            \State $d[v] \gets d[u] + w(u,v)$ \Comment{relax: found a shorter route to $v$}
            \State \textsc{DecreaseKey}$(Q, v, d[v])$; \quad set $\mathrm{pred}[v] \gets u$
        \EndIf
    \EndFor
\EndWhile
\State \Return $d[\cdot]$ (and $\mathrm{pred}[\cdot]$ for path reconstruction)
\end{algorithmic}
\end{algorithm}

Problem 2 asks you to implement the distance-only version using Python's
\texttt{heapq}; Problem 3 adds a $\mathrm{pred}[\cdot]$ array and reconstructs
the actual shortest path to a target by following predecessors backward.

In practice, with a binary heap, the cleanest implementation is the \emph{lazy}
one: instead of \textsc{DecreaseKey}, push a new $(d[v], v)$ pair onto the heap
whenever $d[v]$ improves, and discard a popped pair if its vertex has already
been finalized. This is exactly what the exercise solutions do.

\subsection{Correctness}

\begin{theorem}[Correctness of Dijkstra]
\label{thm:dijkstra}
If all edge weights are non-negative, then when Dijkstra's algorithm adds a
vertex $u$ to $S$, the tentative distance $d[u]$ equals the true shortest-path
distance $\mathrm{dist}(u)$ from $s$ to $u$.
\end{theorem}

\begin{proof}[Proof (greedy stays ahead, by induction on $|S|$)]
We prove the invariant: for every vertex $u$ at the moment it is added to $S$,
$d[u] = \mathrm{dist}(u)$, and moreover $d[u]$ is the length of a genuine
$s$-$u$ path. The base case is $u = s$: $d[s] = 0 = \mathrm{dist}(s)$.

\textbf{Inductive step.} Suppose the invariant holds for all vertices currently
in $S$, and let $u \notin S$ be the next vertex extracted (so $d[u] \le d[v]$
for all $v \notin S$). Because every relaxation sets $d[v]$ to the length of an
actual path, $d[u] \ge \mathrm{dist}(u)$ always; we must show $d[u] \le
\mathrm{dist}(u)$, i.e.\ no path is shorter than $d[u]$.

Consider any $s$-$u$ path $P$. Since $s \in S$ and $u \notin S$, $P$ must leave
$S$ at some point: let $(x,y)$ be the first edge of $P$ with $x \in S$ and $y
\notin S$. By the inductive hypothesis $d[x] = \mathrm{dist}(x)$, and when $x$
was added to $S$ the edge $(x,y)$ was relaxed, so $d[y] \le d[x] + w(x,y)$. The
length of $P$ is at least the length of its prefix up to $y$, which is at least
$\mathrm{dist}(x) + w(x,y) = d[x] + w(x,y) \ge d[y]$. Since $y \notin S$ and $u$
was chosen as the minimum, $d[u] \le d[y]$. Putting it together, the length of
$P$ is at least $d[y] \ge d[u]$. As $P$ was arbitrary, $\mathrm{dist}(u) \ge
d[u]$. Hence $d[u] = \mathrm{dist}(u)$.
\end{proof}

The single place where non-negativity is used is the step ``the length of $P$ is
at least the length of its prefix up to $y$'': this needs every edge \emph{after}
$y$ on $P$ to add a non-negative amount. With a negative edge this can fail, and
the proof -- and the algorithm -- break (see \S\ref{sec:negative}).

\subsection{Running time}

The cost is dominated by the priority-queue operations: $n$ \textsc{ExtractMin}
operations and up to $m$ \textsc{DecreaseKey} (or push) operations. With a
\textbf{binary heap}, each costs $O(\log n)$, giving
\[
    O\big((m + n)\log n\big),
\]
which is $O(m \log n)$ for a connected graph (where $m \ge n - 1$). With a more
sophisticated Fibonacci heap, \textsc{DecreaseKey} is $O(1)$ amortized and the
bound improves to $O(m + n\log n)$, but the binary-heap version is what is used
in practice and in this course.

\subsection{Why Dijkstra fails on negative edges}
\label{sec:negative}

Dijkstra's greedy choice -- ``finalize the nearest unfinished vertex, and never
touch it again'' -- is justified only because, with non-negative weights,
extending a path never decreases its length. A negative edge violates this: a
seemingly-distant vertex might be reachable cheaply via a detour that ends with a
negative edge, but by then Dijkstra has already finalized the relevant vertex.

\begin{example}[A negative edge breaks Dijkstra]
Take vertices $s, a, b$ with directed edges $s \to a$ of weight $2$, $s \to b$
of weight $3$, and $b \to a$ of weight $-2$. The true shortest distance to $a$ is
via $s \to b \to a$: $3 + (-2) = 1$. But Dijkstra extracts $a$ first (tentative
distance $2 < 3$), finalizes $\mathrm{dist}(a) = 2$, and never revisits it -- so
it reports the wrong answer $2$. Problem 6 asks you to construct exactly this
counter-example and confirm that Dijkstra and Bellman--Ford disagree on it.
\end{example}

\section{The Bellman--Ford Algorithm}

When edges may be negative, we need a different approach. The Bellman--Ford
algorithm relaxes \emph{every} edge, $n-1$ times in a row.

\begin{algorithm}[h]
\caption{Bellman--Ford Single-Source Shortest Paths}
\begin{algorithmic}[1]
\State $d[s] \gets 0$; \quad $d[v] \gets \infty$ for all $v \ne s$
\For{$i = 1$ to $n - 1$}
    \For{each edge $(u,v) \in E$}
        \If{$d[u] + w(u,v) < d[v]$}
            \State $d[v] \gets d[u] + w(u,v)$ \Comment{relax}
        \EndIf
    \EndFor
\EndFor
\For{each edge $(u,v) \in E$} \Comment{one extra pass detects a negative cycle}
    \If{$d[u] + w(u,v) < d[v]$}
        \State \Return ``negative cycle reachable''
    \EndIf
\EndFor
\State \Return $d[\cdot]$
\end{algorithmic}
\end{algorithm}

\begin{theorem}[Correctness of Bellman--Ford]
If $G$ has no negative cycle reachable from $s$, then after the $n-1$ rounds of
relaxation, $d[v] = \mathrm{dist}(v)$ for every $v$. Moreover, a further
relaxation is possible on some edge \emph{iff} a negative cycle is reachable from
$s$.
\end{theorem}

\begin{proof}[Proof sketch]
A shortest path that uses no negative cycle has at most $n-1$ edges. By induction
on $i$, after $i$ rounds $d[v]$ is at most the weight of the shortest $s$-$v$
path using at most $i$ edges (relaxing every edge each round can only propagate
the best ``$i$-edge'' distance one more hop). After $n-1$ rounds every
shortest path (at most $n-1$ edges) has been captured. If, in the extra pass,
some edge can still be relaxed, then there is a closed walk on which the total
weight strictly decreases -- i.e.\ a negative cycle; conversely a negative cycle
always admits such a relaxation.
\end{proof}

Each round relaxes all $m$ edges, so Bellman--Ford runs in $O(nm)$ time --
slower than Dijkstra's $O((m+n)\log n)$, the price paid for handling negative
edges and detecting negative cycles. Problem 4 asks you to implement
Bellman--Ford returning \texttt{None} on a negative cycle; Problem 5 verifies
that Dijkstra and Bellman--Ford agree on random \emph{non-negative} graphs.

\section{The Minimum Spanning Tree Problem}

\begin{definition}[Minimum Spanning Tree]
Given a connected, undirected graph $G = (V,E)$ with edge weights $w(e)$, a
\emph{spanning tree} is a subset $T \subseteq E$ of $n-1$ edges that connects all
$n$ vertices and contains no cycle. A \emph{minimum spanning tree} (MST) is a
spanning tree of minimum total weight $\sum_{e \in T} w(e)$.
\end{definition}

An MST is the cheapest way to wire up a network so that every node can reach
every other: think of laying cable to connect cities, or a chip's clock-
distribution network. Two classic greedy algorithms solve it -- Kruskal's and
Prim's -- and both rest on a single structural fact, the cut property.

\subsection{The cut and cycle properties}

\begin{definition}[Cut]
A \emph{cut} of $G$ is a partition of $V$ into two non-empty sets $(S, V
\setminus S)$. An edge \emph{crosses} the cut if it has exactly one endpoint in
$S$.
\end{definition}

To keep the statements clean, we first assume all edge weights are
\emph{distinct} (ties can be broken consistently, e.g.\ by edge index; with that
convention the MST is unique and every claim below goes through).

\begin{theorem}[Cut Property]
\label{thm:cut}
Let $(S, V\setminus S)$ be any cut, and let $e$ be the \emph{minimum-weight} edge
crossing it. Then $e$ belongs to every minimum spanning tree of $G$.
\end{theorem}

\begin{proof}[Proof (exchange argument)]
Let $e = (u,v)$ with $u \in S$, $v \notin S$, and suppose for contradiction that
some MST $T$ does not contain $e$. Adding $e$ to $T$ creates exactly one cycle
$C$ (since $T$ is a spanning tree, $T \cup \{e\}$ has one cycle, which passes
through $e$). The cycle $C$ must cross the cut an even number of times, so it
contains at least one other crossing edge $e' \ne e$. Since $e$ is the
minimum-weight crossing edge and weights are distinct, $w(e) < w(e')$. Now form
$T' = (T \setminus \{e'\}) \cup \{e\}$: removing $e'$ from the cycle keeps the
graph connected and acyclic, so $T'$ is a spanning tree, and $w(T') = w(T) -
w(e') + w(e) < w(T)$. This contradicts the minimality of $T$. Hence every MST
contains $e$.
\end{proof}

\begin{theorem}[Cycle Property]
\label{thm:cycle}
Let $C$ be any cycle in $G$, and let $f$ be the \emph{maximum-weight} edge on
$C$. Then $f$ belongs to \emph{no} minimum spanning tree of $G$.
\end{theorem}

\begin{proof}
Suppose some MST $T$ contains $f = (x,y)$. Removing $f$ splits $T$ into two
components, with $x$ on one side (call its vertex set $S$) and $y$ on the other.
Edge $f$ crosses the cut $(S, V \setminus S)$. The cycle $C$ also crosses this
cut, so besides $f$ it contains another crossing edge $f'$ with (since $f$ is the
unique heaviest on $C$) $w(f') < w(f)$. Replacing $f$ by $f'$ reconnects the two
components into a spanning tree $T'$ with $w(T') < w(T)$, contradicting
minimality.
\end{proof}

The cut property tells us which edges \emph{must} be in the MST; the cycle
property tells us which edges \emph{cannot} be. Problem 10 asks you to verify the
cut property empirically, and Problem 11 uses both properties to classify a given
edge as belonging to \emph{every}, \emph{some}, or \emph{no} MST.

\subsection{Kruskal's algorithm}

Kruskal's algorithm sorts the edges by weight and adds each one that does not
create a cycle -- detected with a union-find (disjoint-set) data structure.

\begin{algorithm}[h]
\caption{Kruskal's Minimum Spanning Tree}
\begin{algorithmic}[1]
\State Sort edges so that $w(e_1) \le w(e_2) \le \cdots \le w(e_m)$
\State $T \gets \emptyset$; \quad initialize union-find with each vertex its own set
\For{$i = 1$ to $m$}
    \State let $e_i = (u, v)$
    \If{$\textsc{Find}(u) \ne \textsc{Find}(v)$} \Comment{$u, v$ in different components}
        \State $T \gets T \cup \{e_i\}$; \quad $\textsc{Union}(u, v)$
    \EndIf
\EndFor
\State \Return $T$
\end{algorithmic}
\end{algorithm}

\begin{theorem}
Kruskal's algorithm produces a minimum spanning tree.
\end{theorem}

\begin{proof}
Every edge $e = (u,v)$ that Kruskal \emph{adds} is the minimum-weight edge
crossing the cut $(S, V \setminus S)$, where $S$ is the set of vertices in $u$'s
current component: $e$ connects two different components, and every cheaper edge
was already processed and either added (staying inside a component) or
discarded -- in particular, no cheaper edge crosses this cut, since had one
existed it would have merged the components earlier. By the cut property
(Theorem~\ref{thm:cut}) $e$ belongs to every MST, so every edge Kruskal adds is
safe. Since Kruskal adds edges until all vertices are connected (it never creates
a cycle), it returns a spanning tree, all of whose edges belong to the MST --
hence the MST itself.
\end{proof}

\paragraph{Union-find and running time.} The bottleneck is sorting the edges:
$O(m \log m) = O(m \log n)$. The union-find operations -- with \emph{path
compression} and \emph{union by rank} -- cost nearly $O(1)$ amortized (formally,
$O(\alpha(n))$ where $\alpha$ is the inverse Ackermann function, effectively
constant), so the total is $O(m \log n)$. The \texttt{UnionFind} class in
\texttt{starter\_code.py} implements both optimizations; Problem 7 asks you to
build Kruskal's algorithm on top of it and return both the MST edge set and its
total weight.

\subsection{Prim's algorithm}

Prim's algorithm grows a single tree outward from an arbitrary start vertex,
always adding the cheapest edge that leaves the current tree -- the direct graph
analogue of Dijkstra, using a heap of candidate crossing edges.

\begin{algorithm}[h]
\caption{Prim's Minimum Spanning Tree}
\begin{algorithmic}[1]
\State Pick any start vertex $r$; \quad $S \gets \{r\}$; \quad $T \gets \emptyset$
\State Push every edge incident to $r$ onto a min-heap $Q$, keyed by weight
\While{$|S| < n$ and $Q \ne \emptyset$}
    \State $(w, v) \gets \textsc{ExtractMin}(Q)$ \Comment{cheapest candidate edge}
    \If{$v \in S$} \textbf{continue} \Comment{stale entry; skip}
    \EndIf
    \State $S \gets S \cup \{v\}$; \quad add the chosen edge to $T$
    \For{each edge $(v, x)$ with $x \notin S$}
        \State Push $(w(v,x), x)$ onto $Q$
    \EndFor
\EndWhile
\State \Return $T$
\end{algorithmic}
\end{algorithm}

\begin{theorem}
Prim's algorithm produces a minimum spanning tree.
\end{theorem}

\begin{proof}
At every step, $S$ is the set of vertices in the current tree, and Prim adds the
minimum-weight edge crossing the cut $(S, V \setminus S)$. By the cut property
(Theorem~\ref{thm:cut}) this edge belongs to every MST, so every edge added is
safe. After $n-1$ such additions, $S = V$ and $T$ is a spanning tree composed
entirely of MST edges -- the MST.
\end{proof}

Using a binary heap, Prim's algorithm runs in $O((m+n)\log n)$, the same as
Dijkstra. Problem 8 asks you to implement it (returning the total MST weight),
and Problem 9 verifies that Kruskal and Prim always agree on the total weight
for random connected graphs -- a useful empirical check, since (with ties) they
may pick different edge sets but must achieve the same minimum total.

\section{Application: Single-Link Clustering}

A pleasant payoff of the MST machinery is an optimal clustering algorithm.
Suppose we have $n$ points (objects) with a pairwise distance between each pair,
and we want to partition them into $k$ groups (\emph{clusters}). A natural
objective is to \emph{maximize the spacing}: the \emph{spacing} of a clustering
is the minimum distance between any two points that lie in \emph{different}
clusters, and we want this minimum gap to be as large as possible.

\begin{definition}[$k$-Clustering of Maximum Spacing]
Given a complete weighted graph on the points (weights = distances) and an
integer $k$, partition the vertices into $k$ non-empty clusters so that the
minimum weight of an edge joining two different clusters (the \emph{spacing}) is
maximized.
\end{definition}

\begin{theorem}[Single-link clustering via Kruskal]
The following greedy procedure produces a $k$-clustering of maximum spacing: run
Kruskal's algorithm on the points, but \emph{stop} as soon as the union-find has
exactly $k$ components. Equivalently, build the MST and delete its $k-1$
most-expensive edges; the remaining $k$ components are the clusters, and the
spacing equals the weight of the $(k-1)$-th most expensive MST edge (the next
edge Kruskal would have added).
\end{theorem}

\begin{proof}[Proof sketch]
Running Kruskal until $k$ components remain is exactly adding the $n-k$ cheapest
``component-merging'' edges; the spacing of the resulting clustering is the
weight $d^\ast$ of the next edge Kruskal would add. Consider any other
$k$-clustering $\mathcal{C}'$ that differs from ours. Then two points $p, q$ in
the same cluster of ours lie in different clusters of $\mathcal{C}'$. The Kruskal
path joining $p$ and $q$ uses only edges of weight $\le d^\ast$, and somewhere
along it two consecutive points are split by $\mathcal{C}'$ -- giving
$\mathcal{C}'$ a crossing edge of weight $\le d^\ast$, so its spacing is at most
$d^\ast$. Hence no clustering beats ours.
\end{proof}

Problem 12 asks you to implement this: run Kruskal until $k$ components remain
and return both the clustering and the spacing.

\section{Summary and Looking Ahead}

This week applied greedy design to two cornerstone graph problems:

\begin{itemize}
    \item \textbf{Shortest paths.} \emph{Dijkstra's algorithm} greedily
    finalizes the nearest unprocessed vertex; the ``greedy stays ahead'' proof
    (Theorem~\ref{thm:dijkstra}) works only for non-negative weights, in
    $O((m+n)\log n)$ with a binary heap. A single negative edge breaks it,
    motivating \emph{Bellman--Ford} ($O(nm)$), which also detects negative
    cycles.
    \item \textbf{Minimum spanning trees.} The \emph{cut property} and
    \emph{cycle property} (Theorems~\ref{thm:cut}, \ref{thm:cycle}) determine
    exactly which edges must, and must not, appear in an MST. Both
    \emph{Kruskal's} algorithm (sort + union-find) and \emph{Prim's} algorithm
    (grow + heap) are correct because every edge they add is a minimum crossing
    edge of some cut.
    \item \textbf{Application.} Single-link $k$-clustering is just Kruskal
    stopped early, and it provably maximizes the spacing.
\end{itemize}

The unifying lesson is that one structural fact (the cut property) justifies two
different-looking greedy algorithms, just as one proof template (greedy stays
ahead) justified Dijkstra. Next week we change paradigms entirely: from greedy
to \textbf{divide and conquer} (Kleinberg \& Tardos Ch.\ 5), where we solve a
problem by splitting it into independent subproblems, solving each recursively,
and combining -- with the running time governed by recurrence relations and the
Master Theorem.

\end{document}
